\section{Giới thiệu}
\label{sec:intro}

\subsection{Bối cảnh}

Retrieval-Augmented Generation (RAG)~\cite{lewis2020rag} là cách phổ biến nhất hiện nay để đưa tri
thức bên ngoài vào mô hình ngôn ngữ lớn (LLM): với mỗi câu hỏi, một bộ truy hồi (retriever) tìm
$k$ đoạn văn liên quan trong kho tài liệu, rồi LLM sinh câu trả lời dựa trên các đoạn đó. Để không bỏ
sót bằng chứng, hệ thống thường lấy $k$ khá lớn (10--20 đoạn). Điều này gây ra ba vấn đề:

\begin{enumerate}
  \item \textbf{Chi phí và độ trễ.} Prompt dài hàng nghìn token; chi phí API và thời gian xử lý
  prompt (prefill) tăng tuyến tính (thậm chí bậc hai với attention) theo độ dài ngữ cảnh.
  \item \textbf{Nhiễu.} Phần lớn các đoạn truy hồi không chứa đáp án; chúng có thể đánh lạc hướng LLM.
  \item \textbf{Thiên lệch vị trí.} \citet{liu2024lost} chỉ ra hiện tượng \emph{lost in the middle}:
  độ chính xác của LLM cao khi đoạn chứa đáp án nằm ở đầu hoặc cuối prompt, nhưng giảm mạnh khi nó
  nằm ở giữa (đường cong hình chữ U).
\end{enumerate}

Hướng \emph{Efficient RAG} giải quyết các vấn đề trên ở hai đầu của pipeline: (i) \textbf{nén ngữ
cảnh} (context/prompt compression) -- loại bỏ token, câu hoặc đoạn không cần thiết trước khi đưa vào
LLM; và (ii) \textbf{truy hồi hiệu quả} (efficient retrieval) -- giảm thời gian, bộ nhớ và chi phí
của bước tìm kiếm mà vẫn giữ chất lượng.

\subsection{Bài toán}

Cho câu hỏi $q$ và tập $n=20$ đoạn văn truy hồi được $D=(d_1,\dots,d_n)$, trong đó đúng một đoạn
$d_g$ chứa đáp án (đoạn ``vàng''). Một bộ nén $C$ biến $(q, D)$ thành ngữ cảnh ngắn hơn
$c = C(q, D)$, sau đó LLM sinh câu trả lời $\hat a = \mathrm{LLM}(q, c)$. Mục tiêu là đồng thời:
\begin{itemize}
  \item tối đa độ chính xác $\mathrm{Acc} = \Pr[\hat a \text{ chứa đáp án đúng}]$;
  \item tối thiểu số token $|c|$ (tức tối đa tỉ lệ nén $|D|/|c|$), độ trễ và chi phí;
  \item ổn định khi vị trí $g$ của đoạn vàng thay đổi (bền vững với thiên lệch vị trí).
\end{itemize}
Ở phía truy hồi, cho kho $\mathcal{C}$ gồm khoảng 50 nghìn đoạn văn, ta so sánh các bộ truy hồi theo
chất lượng (tỉ lệ top-$k$ chứa đáp án) và hiệu năng (kích thước chỉ mục, thời gian truy vấn).

\subsection{Câu hỏi nghiên cứu và đóng góp}

Báo cáo trả lời năm câu hỏi:
\begin{description}
  \item[RQ1.] Các bộ nén ngữ cảnh từ thế hệ đầu (LLMLingua-2, LongLLMLingua, 2023--2024) đến thế hệ
  \emph{pruner} mới (Provence -- ICLR 2025, OpenProvence -- 2025) giảm được bao nhiêu token mà không
  làm giảm độ chính xác? Chúng có làm ``phẳng'' đường cong chữ U không?
  \item[RQ2.] Các kỹ thuật truy hồi hiệu quả (BM25S, static embedding, Matryoshka, lượng tử hoá
  int8/binary, HNSW) đánh đổi chất lượng -- tốc độ -- bộ nhớ thế nào?
  \item[RQ3.] Kết quả của nhóm so với các công bố trên cùng bộ dữ liệu ra sao?
  \item[RQ4.] Không có GPU thì pipeline nén + truy hồi có còn dùng được không?
  \item[RQ5.] Nén mềm \emph{encode-once} (mỗi đoạn được mã hoá một lần, offline, thành vài vector) có thay
  được văn bản khi đọc bằng LLM nhỏ không?
\end{description}

Đóng góp chính của nhóm:
\begin{itemize}
  \item Xây dựng thư viện mã nguồn mở \texttt{effrag} (Python, \texttt{uv}) với giao diện thống nhất
  cho 10 phương pháp nén (8 dạng văn bản, 2 nén mềm) và 7 cấu hình truy hồi, cùng giao thức đánh giá tái lập được (seed cố định,
  bộ đếm token chung, cache kết quả LLM).
  \item Đánh giá có hệ thống trên bộ NQ-Open 20 tài liệu của \citet{liu2024lost} -- bộ dữ liệu mà
  LongLLMLingua~\cite{jiang2024longllmlingua} và LLMLingua-2~\cite{pan2024llmlingua2} đã công bố kết
  quả theo từng vị trí. Theo khảo sát của nhóm, chưa có công bố nào đánh giá họ pruner Provence trên
  đúng giao thức này.
  \item Đánh giá nén mềm trên cùng giao thức: PISCO~\cite{louis2025pisco} (Mistral-7B) và một bộ nén
  encode-once do nhóm tự huấn luyện trên Qwen2.5-1.5B, có biến thể hybrid giữ đoạn quan trọng nhất ở dạng
  văn bản; với một cổng go/no-go đặt trước và chẩn đoán chỉ ra vì sao mô hình nhỏ chưa thay được văn bản.
  \item Một pipeline RAG end-to-end (retriever lượng tử int8 $\to$ bộ nén $\to$ reader) trên kho
  50,7k đoạn, xác nhận các kết luận của bộ dữ liệu có kiểm soát trong điều kiện truy hồi thật.
  \item Phân tích đánh đổi độ chính xác -- token -- độ trễ và một demo tương tác (Gradio) minh hoạ
  trực quan phần ngữ cảnh mà mỗi phương pháp giữ lại.
\end{itemize}

\subsection{Quy ước thuật ngữ}
\label{sec:terms}

Báo cáo giữ nguyên tiếng Anh cho các thuật ngữ đã thành tên gọi chuẩn trong tài liệu về RAG và khó
dịch mà không mất nghĩa (reader, retriever, rerank, embedding, token, prompt, accuracy, context
recall). Các khái niệm phổ thông dùng tiếng Việt (nén ngữ cảnh, truy hồi, chỉ mục, lượng tử hoá, độ
trễ). Bảng~\ref{tab:glossary} nêu nghĩa của các thuật ngữ dùng xuyên suốt báo cáo.

\begin{table}[!ht]
\centering
\caption{Thuật ngữ và ký hiệu dùng trong báo cáo.}
\label{tab:glossary}
\small
\renewcommand{\arraystretch}{1.12}
\begin{tabularx}{\linewidth}{@{}>{\raggedright\arraybackslash}p{0.27\linewidth} X@{}}
\toprule
Thuật ngữ & Nghĩa trong báo cáo \\
\midrule
Reader & LLM đọc ngữ cảnh và sinh câu trả lời (thí nghiệm chính: Qwen2.5-7B-Instruct). \\
Retriever, rerank & Retriever tìm các đoạn liên quan trong kho; rerank chấm điểm lại các đoạn đã có
  bằng một cross-encoder (reranker) để giữ $k$ đoạn tốt nhất. \\
Dense retriever, embedding & Truy hồi bằng vector embedding của câu hỏi và đoạn văn (bge-small),
  khác BM25 vốn dựa trên trùng khớp từ vựng. \\
Đoạn vàng (gold passage) & Đoạn duy nhất chứa đáp án trong 20 đoạn; vị trí của nó (1, 10, 20) là
  biến của thí nghiệm. \\
Full, Closed-book, Oracle & Đưa đủ 20 đoạn / không đưa đoạn nào / chỉ đưa đoạn vàng vào prompt. \\
Accuracy & Tỉ lệ câu hỏi có output chứa một đáp án chuẩn (best-subspan EM của \citet{liu2024lost}). \\
Context recall & Tỉ lệ câu hỏi mà đáp án còn nằm trong ngữ cảnh sau nén; đo được không cần LLM. \\
Tỉ lệ nén & Số token ngữ cảnh gốc chia cho số token sau nén (tiktoken \texttt{o200k\_base}). \\
Nén cứng / nén mềm & Hard / soft compression: đầu ra là văn bản (dùng được với mọi LLM) / là vector
  đưa thẳng vào LLM. \\
Pruning & Cắt bỏ phần không liên quan của ngữ cảnh ở mức câu hoặc đoạn (Provence). \\
Vector nhớ (memory slot) & Vector do bộ nén mềm sinh ra để thay cho nội dung một đoạn. \\
Encode-once, hybrid & Mỗi đoạn được mã hoá một lần, offline, độc lập với câu hỏi; biến thể hybrid
  giữ đoạn xếp hạng cao nhất ở dạng văn bản, các đoạn còn lại ở dạng vector nhớ. \\
Cổng go/no-go & Điều kiện đặt trước trên tập dev: đạt thì đánh giá tiếp, không đạt thì dừng và
  phân tích nguyên nhân. \\
Lost in the middle & Hiện tượng LLM dùng kém thông tin nằm giữa ngữ cảnh (đường cong chữ U theo vị
  trí đoạn vàng). \\
Paired bootstrap, CI & Kiểm định chênh lệch giữa hai cấu hình trên cùng tập câu hỏi; CI là khoảng tin
  cậy 95\%. ``Có ý nghĩa'' nghĩa là CI không chứa 0. \\
Lượng tử hoá int8 / binary & Lưu mỗi chiều của embedding bằng 8 bit / 1 bit thay vì 32 bit. \\
\bottomrule
\end{tabularx}
\end{table}

Phần còn lại của báo cáo: Mục~\ref{sec:related} tổng quan nghiên cứu liên quan; Mục~\ref{sec:method}
mô tả phương pháp; Mục~\ref{sec:impl} trình bày cài đặt; Mục~\ref{sec:results} phân tích kết quả;
Mục~\ref{sec:conclusion} kết luận.
